\chapter{A Difference Appears}

\section{The first fact}

Before a thing can be measured it must first be told apart from what it is not, and that telling-apart is the whole
of the instrument's ground floor. Nothing here is counted, weighed, or valued. There is only a \emph{this} that is
not a \emph{that} --- the smallest possible act, prior to number, on which every later act is built. This first
section builds that act and only that act, and it is worth building slowly, because everything the machine will ever
do is this one step, iterated.

In this volume the first fact is read as a piece of type theory, and the reading pays off at once. A difference here
is not an event that happens and is then recorded; it is a \emph{type former}. To tell two things apart is to exhibit
an \emph{inhabitant} of a type whose very existence \emph{is} the distinction --- a term the machine constructs, not
a fact it files away. This is the propositions-as-types reading, and it is the register the whole book stands in: a
proposition simply \emph{is} the type of its proofs, so to hold the proposition ``these two differ'' is to hold a term
of that type, and constructing the term is the entire act. We never assert a difference and then go hunting for
evidence. Under this reading the difference \emph{is} the evidence.

Three steps carry the ground floor. The first is \emph{formation}, and it comes before everything. The distinguishable
is introduced by a formation rule: given two carriers, it forms a type --- the type of ways the two could be told
apart. This is the thinnest former there is; it says only \emph{what it would take} to separate the two, not that they
are in fact separated. The ordering matters, and the book leans on it throughout: formation (the type exists) precedes
introduction (an inhabitant is built) precedes elimination (the distinction is finally \emph{used}, downstream, when
the second chapter counts). Measurement lives three steps up from here. This section is pure formation --- the type,
and nothing yet put in it.

The second step says what the ``level'' is that does the telling-apart, and the answer is the surprising one. The
single running index the machine carries is not a counter it invented; it is the \emph{universe hierarchy} of
dependent type theory --- the tower of \(\mathsf{Type}\) upon \(\mathsf{Type}\) that Martin-L\"of~\cite{martin1970}
built --- pressed into one humble service: to give the word ``different'' a mechanical, \emph{decidable} meaning, so
that a machine and not a person settles it. Two things are distinguishable precisely when they are separated in that
hierarchy, and the comparison of levels is a primitive judgment the machine can simply make --- not an arithmetic
performed on integers, but a decision read straight off the typing. This is Leibniz's~\cite{leibniz1686} identity of
indiscernibles run backwards: he held that things sharing every property are one; the machine holds that two things
are one until some level tells them apart, and ``indiscernible'' means here, natively, that no term separates them.
The label the machine reads is the one it already tracks at the universe level, and it carries that label without a
single integer --- the level is a typing, not a number.

The third step is that this costs no allocation, and in this register that is not a claim about speed but about
\emph{where the equation lives}. The distinction is not written into a cell and fetched back later to be compared; it
is available \emph{definitionally}, by the same judgmental machinery that already tracks levels to keep the machine's
descriptions coherent. Nothing is stored, because nothing needs to be looked up. Sameness here is the sameness of
reflexivity --- decided by computation on the spot, never retrieved from a memory. The ledger, in other words, was
already open before the first fact was written; the fact did not have to buy itself a page.

And now the first of the book's two working verbs, introduced where it is native. To \emph{tange} is to select --- to
pick a \emph{this} out of a field by the characteristic that parts it from its \emph{that}. Telling two things apart
is not yet counting them; it is the pure act of selection, and that act is a tange. The word earns its keep exactly
here, because to inhabit the distinguishable \emph{is} to select one carrier against another by the level at which
they part: the tange \emph{is} the inhabitant. Its partner verb --- to \emph{funge}, to place a thing in a bag of its
likes --- is not yet available, and cannot be, for you cannot bag things by their sameness until you have first tanged
them apart. That comes in the next chapter, where counting begins: to count is to funge, to bag distinguishables
already tanged into a single like heap whose size is the number. So the order of the two verbs is set here at the
start, and it is the order of the whole construction: tange precedes funge, as selection precedes bagging, as
formation precedes elimination.

This lands the reader on the one fact the ground floor exists to establish: \emph{resolution is logically upstream of
measurement}. You must have a term before you can have a number of terms; inhabitation comes before quantity. It is
the same point Nyquist~\cite{nyquist1928} and Shannon~\cite{shannon1948} would later make about signals --- that you
cannot read past a resolution your instrument cannot resolve --- made here one level lower, about symbols: you cannot
read a scale whose marks you cannot first tell apart into distinct terms. The mark must be a \emph{this-not-that}
before it can be a \emph{how-many}.

It is worth marking, at the last, how much of the book is already latent in this one step. Everything the instrument
will later call a number, a charge, a mass, a cost is built by iterating exactly this: take a difference, let it be
its own label, and pay only for the turn. The naming of the electron eight chapters on, and the honest bracket the
machine hands back at the very end, are this same act matured --- not new machinery, but the first fact taken around
the loop enough times to be read. And notice what has \emph{not} appeared: no number, not a count, not a magnitude.
The first fact is pre-numeric by its very construction, which makes this the cleanest place in the book to adopt the
discipline the whole instrument keeps --- it will not reach for a magnitude it cannot earn, and the reason is visible
already here, at formation, where there is not yet even a number to reach for.

\section{The symbol as a label}

The first fact was a difference; the natural next question is where the machine \emph{puts} it. Having told two things
apart, where does it write the distinction down, so that it can find it again? The answer is the hinge of this chapter,
and it is a small shock: nowhere. The machine writes the difference nowhere, because the difference is already its own
record. The distinction and the label of the distinction are not two things joined by an act of naming; they are one
object. And because they are one object, the label costs nothing to keep. Difference, in this machine, is free.

To see why, ask what the inhabitant of the last section was actually \emph{made of}. In the ordinary account there are
two separate things with a step between them: a difference on one side, a name for the difference stored somewhere on
the other, and a copy that carries the first into the second. Here that middle step is missing, and its absence is the
whole content of the section. The label the machine hands you is exactly the universe level at which the two carriers
parted --- the very type at which they differ. Telling apart already \emph{produced} that type; there is no second act
in which the machine turns around and records what it found. The tange of the last section did not select a difference
and then go write its name; the selection \emph{was} the name.

This is why the claim ``the difference is its own label'' is not a play on words. In this register a label just
\emph{is} a position in the universe hierarchy --- which level, relative to what --- and the tange delivered exactly
such a position. To select a \emph{this} against a \emph{that} by the level at which they part is already to hold that
level; and the level is the label. Naming adds nothing, because the tange already named. The characteristic you
selected by and the record you walk away with are the same object, and the reader who was holding them apart as two
may now let them fall together into one.

The consequence is that the distinction requires no allocation of its own --- but the claim has to be scoped exactly,
because scoped loosely it is false. The machine does not reserve storage \emph{for the difference}; it reads the
distinction straight off the one index it must track anyway to keep its own descriptions coherent, the compiler's
universe level. It does not follow that the machine reserves \emph{nothing}. The construction carries one standing
datum, a shared truth-constant, and it is honest about it in its own voice: \emph{``we are grabbing the label the
compiler uses to track universe level. There is no memory allocation so far, other than \texttt{Fact.Truth}, which is
on the DATA page, not the INSTRUCTION page.''} Read the carve-out as carefully as the claim. What is free is the
\emph{difference} --- nothing is set aside per distinction, because the label lives at a level the machine already
keeps --- while the one thing that is allocated, the shared truth-constant, is a single data-page fixture, not a fresh
cell spent on each telling-apart. ``Free'' here does not mean fast, and it does not mean cheap; those are claims about
a clock, and this is not a claim about a clock. It means \emph{nothing added}: the difference is decided where it
stands, not stored to be looked up --- free in the way a definitional equality is free.

Underneath all of this is an idea far older than any computer: a symbol means nothing in isolation, and means only by
its contrast with the symbols it is not. Value is difference --- Saussure~\cite{saussure1916} said it of words, and
the theory of information later made it quantitative, a symbol carrying meaning exactly to the degree it could have
been another. This machine takes the old idea at its word and builds on it literally: its symbols \emph{are} their
differences, a symbol defined as its inequality to every other. The label is the contrast itself, not a name pinned on
after the contrast was noticed. The semiotics is illumination here, texture around the fact; what the machine actually
carries is the inequality, and nothing of Saussure's beyond it.

So take stock of what the instrument has handed you at the end of the ground floor: not a measurement, and not yet a
number --- a \emph{position}. Which level, relative to what. That is all of it, and it is exactly enough, because a
position told apart from other positions is the seed of everything the machine will grow. And the label is about to
earn its keep a second way: in the very next chapter the machine will \emph{funge} --- gather distinguishables that
share a label into a single bag --- and the size of that bag will be its first number. The label is precisely what the
count will bag by; but that is the next chapter's work, and it performs no count here. What this section settles is
only that the step onto the count carries no hidden storage to be paid for later. Everything the machine will grow ---
number, and the rest --- is this one act iterated: take a difference, let it be its own label, and pay only for the
turn. The only bill it ever runs up is the turn, never the label; and the turns are the whole subject of the chapters
to come.

\section{What a difference is not}

The ground floor is nearly built, and the last thing to do on it is the most easily skipped: to say what the first
fact is \emph{not}. A difference is not yet a number. It is not a measurement, not a magnitude, not a count --- the
smallest possible act, and no more. The temptation, having built something, is to read more into it than it holds; the
whole discipline of this book is the refusal of that temptation, and the refusal is easiest to learn here, where there
is least to inflate.

The relation the machine actually has is inequality, and inequality alone. To ask whether two things differ is to ask
whether the one is unequal to the other, and the answer comes back \emph{same} or \emph{not-same} and stops.
Inequality is the coarsest relation two things can stand in --- it reports a bare yes-or-no and carries nothing else
along with it. It is worth pausing on how much that ``nothing else'' rules out, because each thing it rules out is a
structure the reader may be tempted to assume is already present. There are three, and each is a refusal.

The first refusal: a difference is \emph{not an order}. To be a number is, at the very least, to sit in a chain --- to
be comparable, to have a \emph{less-than} and a \emph{greater-than}, an order written \(\le\). The distinguishable
tier exposes no such thing. Its one relation is symmetric and blind to direction: it can report that two things are
not the same, but not which of them comes first, or is larger, or sits higher. Order does not live on the distinction
at all; it is built higher up, on the numeric objects the later tiers construct --- never as a feature of the bare
difference. A reader who searches the ground floor for a \emph{which-is-bigger} will not find one, and the reason is
not that it is hidden but that the question cannot yet be posed. At this level, ``which is bigger?'' does not so much
as parse.

The second refusal: a difference is \emph{not a magnitude}. A measurement answers \emph{how much}; a difference
answers only \emph{whether}. There is no unit here, no scale, no distance --- nothing that would let the machine say
two things are far apart or close, only that they are two. The datum is nominal in the plainest sense: same or
different, and no finer. (One might recall Frege's~\cite{frege1884} remark that a number belongs to a concept and not
to a bare object --- but only in passing, for the machine's refusal here is not a philosopher's; it is structural, a
matter of which fields the tier does and does not carry.)

The third refusal is the one this book is built to make precise, and it is cleanest stated in the two working verbs.
A difference is a pure \emph{tange} --- a selection, a \emph{this} picked out against a \emph{that}. A count is
something else entirely: it is the size of a \emph{funge}, the cardinality of a bag of like things. And you cannot
take the size of a bag before you have a bag, nor bag things by their likeness before you have first tanged them
apart. So ``not yet a number'' is not a shortfall the machine will outgrow by trying harder; it is an \emph{ordering},
as strict as any in the construction. The tange comes first; the funge comes after; and number lives entirely on the
funge side of that line. A number is a \emph{later construction} --- something the next chapter builds by bagging
distinguishables the ground floor has already tanged apart --- not a reading available here at formation. The count
sits at a strictly higher tier than the distinction, and no amount of staring at a single difference will yield one.

That, in the end, is the shape of the whole discipline, seen in miniature at the smallest possible scale. The
machine's refusal to call a position a quantity is not timidity, and it is not incompleteness; it is correctness. It
holds exactly one thing --- a distinction, a position, which level relative to what --- and it declines to inflate
that one thing into a number it has not yet built. This is the same refusal, grown up, that the instrument will make
at the very end, when it declines to call its honest bracket the real value the bracket contains: in both places the
machine reports exactly what it has and not a hair more. The ground floor and the final floor keep the same
discipline; what stands between them is only the long climb --- and the climb begins in the next chapter, when the
machine takes the differences it has tanged apart and, for the first time, begins to funge them into a count.
